% Generated draft: requires mathematical and source review.
\documentclass[11pt]{article}
\usepackage[margin=25mm]{geometry}
\usepackage{fontspec}
\setmainfont{Latin Modern Roman}
\usepackage{amsmath,amssymb,mathtools,unicode-math}
\setmathfont{Latin Modern Math}
\usepackage{enumitem,needspace,xurl,hyperref}
\setlength{\parindent}{0pt}
\setlength{\parskip}{5pt}
\setlength{\emergencystretch}{3em}
\newcommand{\N}{\mathbb{N}}
\newcommand{\Z}{\mathbb{Z}}
\newcommand{\Q}{\mathbb{Q}}
\newcommand{\R}{\mathbb{R}}
\newcommand{\C}{\mathbb{C}}
\newcommand{\F}{\mathbb{F}}
\newcommand{\A}{\mathbb{A}}
\newcommand{\PP}{\mathbb{P}}
\newcommand{\E}{\mathbb{E}}
\newcommand{\eps}{\varepsilon}
\newcommand{\dd}{\,\mathrm{d}}
\begin{document}
\section*{UnsolvedMath ID: 30002124}
\textbf{Source problem code:} catalogue entry
\subsection{Definitions and mathematical statement}
The supplied record bundles three independent questions. They are stated separately below, with unresolved hypotheses and conventions made explicit.

\begin{enumerate}
\item \textbf{Discrete Green and kernel estimates.} Let \(\Gamma\) and \(\Gamma^*\) be planar embedded graphs that are dual to one another, with corresponding primal and dual edges meeting orthogonally. The source asks whether, under unspecified ``mild regularity assumptions,'' one has
\[
G(u;u_0)=\frac{1}{2\pi}\log|u-u_0|+O(1),
\qquad
K(v;z_0)=O\bigl(|v-z_0|^{-1}\bigr).
\]
Here the absolute values denote distances between the relevant embedded points. The record does not define the functions \(G\) and \(K\), their arguments' domains, the discrete operators and weights involved, or their normalizations and boundary conditions. In particular, the designation of \(G\) as a Green function does not determine its sign or additive normalization.

For specified domains and an asymptotic regime, the requested uniform bounds would mean estimates of the form
\[
\left|G(u;u_0)-\frac{1}{2\pi}\log|u-u_0|\right|\leq C_G,
\qquad
|K(v;z_0)|\leq\frac{C_K}{|v-z_0|},
\]
for the intended distinct pairs of arguments, with constants independent of the variables over which uniformity is claimed. The admissible pairs, limiting regime, regularity hypotheses, and permitted dependence of these constants are not supplied. Thus these inequalities describe the requested estimates but do not yet constitute a fully specified conjecture.

\item \textbf{Escaping geodesic rays on a cone surface.} A Euclidean cone surface is a surface with an intrinsic metric locally Euclidean away from cone singularities and locally modeled at each cone singularity on a Euclidean cone. A cone point of total angle \(\theta\) has angle defect \(2\pi-\theta\); negative curvature at that point means \(\theta>2\pi\). Let \(S\) be a metrically complete, noncompact Euclidean cone surface, and let \(\Sigma\) be its set of cone singularities. Assume, as stated in the record, that it has finitely many negatively curved cone singularities.

The question asks whether every point \(x\in S\) is the initial point of a geodesic ray \(\gamma:[0,\infty)\to S\) that avoids every singularity and escapes to infinity. Escape to infinity means
\[
\forall\, C\subseteq S\text{ compact}\quad
\exists\,T\geq0\quad
\forall\,t\geq T:\ \gamma(t)\notin C.
\]
The avoidance condition requires clarification: if it means
\(\gamma([0,\infty))\cap\Sigma=\varnothing\), then it is incompatible with \(\gamma(0)=x\) for singular initial points. Two possible formulations are to restrict the initial point to \(S\setminus\Sigma\), or to allow any initial point and require only
\(\gamma((0,\infty))\cap\Sigma=\varnothing\). The record does not choose between them.

The meaning of ``geodesic ray'' also needs specification: a globally minimizing unit-speed ray satisfies
\[
d_S(\gamma(s),\gamma(t))=|s-t|\qquad(s,t\geq0),
\]
whereas a unit-speed locally geodesic curve need not satisfy this identity globally. The source does not distinguish these meanings. It also does not say whether all singularities are the finitely many negatively curved ones, or whether additional singularities of other curvature signs are allowed, nor does it specify boundary conventions. These qualifications remain unresolved.

\item \textbf{Diameter after gluing random polygon triangulations.} For an integer \(n\geq3\), let \(\mathcal T_n\) denote the intended finite set of triangulations of an \(n\)-gon. Choose \(T_n^{(1)}\) and \(T_n^{(2)}\) independently and uniformly from this set, so that
\[
\mathbb P\bigl(T_n^{(1)}=t_1,\ T_n^{(2)}=t_2\bigr)
=|\mathcal T_n|^{-2}
\qquad(t_1,t_2\in\mathcal T_n).
\]
Glue the two triangulations using the intended boundary identification, and let \(H_n\) be the resulting graph. For vertices \(a,b\), let \(d_{H_n}(a,b)\) be the minimum number of edges in a path joining them, and define its graph diameter by
\[
D_n=\max_{a,b\in V(H_n)}d_{H_n}(a,b).
\]
The question is to determine the typical or expected order of \(D_n\); the supplied conjectural exponent is \(1/4\).

If ``typical'' is intended to mean concentration of the exponent in probability, a precise candidate formulation is
\[
\forall\,\varepsilon>0,\qquad
\lim_{n\to\infty}
\mathbb P\!\left(n^{1/4-\varepsilon}\leq D_n\leq n^{1/4+\varepsilon}\right)=1.
\]
If the intended target is instead the expectation, the candidate formulation is
\[
\mathbb E[D_n]=n^{1/4+o(1)}\qquad(n\to\infty).
\]
These are distinct targets, and the source does not specify which is conjectured. It also leaves the precise triangulation ensemble, labeling or rooting conventions, boundary gluing rule, and treatment of identified edges unspecified. The probability and expectation formulations above are therefore explicitly conditional interpretations, not additional source assertions.
\end{enumerate}
\subsection{Short English statement}
The record contains three independent questions: uniform estimates for discrete kernels on orthogonal dual planar graphs, existence of singularity-avoiding escaping geodesic rays on complete noncompact Euclidean cone surfaces, and the diameter of two glued independent uniform polygon triangulations, with conjectural exponent \(1/4\). Essential operator definitions, geometric conventions, and probabilistic specifications are missing. The supplied open status and verification labels require independent review.
\subsection{Sources}
\begin{itemize}
\item Supplied local UnsolvedMath registry record, permanent ID \(30002124\), titled ``Discrete Green Kernels, Escaping Geodesics, and Random Triangulations''; exact source problem code \texttt{catalogue entry}. Collection: Oberwolfach Reports Open Problems. The supplied proposed year is 2012; no individual proposer is identified.
\item Supplied report citation: \emph{Discrete Differential Geometry (2013)}, Oberwolfach Reports, pp.\ 2077--2137, DOI \texttt{10.4171/owr/2012/34}. Supplied source link: \url{https://doi.org/10.4171/owr/2012/34}. The citation's 2013 designation and the record's 2012 year information are retained without resolving their relationship. No question-specific page or speaker attribution is supplied.
\item Supplied literature source, labeled ``OWR report'': \url{https://oa.tib.eu/renate/server/api/core/bitstreams/76ac1159-94b4-4a15-9fc6-21b50a98dcc0/content}.
\item Literature assessment dated exactly \texttt{2026-08-21}: the supplied curation labels the record open, describes it as incomplete and conflating a geodesic-ray avoidance question with a random-triangulation diameter question, and says that related results are partial and the combined extracted claim remains unresolved. This is supplied dated curation, not an independently verified literature assessment.
\item The supplied statement status is \texttt{corrected\_verified}. Its accompanying verification text says that the bundled extraction was retained as three explicitly separated questions because all three occur as independent open problems in the same official report. This is a supplied verification claim; neither the report nor the current mathematical status has been independently checked for this draft.
\end{itemize}
\end{document}
